\subsection{Limits of Gamma-normal closure and repeated reflection}
\label{sec:selected-gamma-limits}

The selected faithful free factor does not by itself give faithfulness
of its \(\Gamma\)-stable closed normal closure. The precise quotient
and the examples below explain this limitation. All group and module
examples in this subsection are formal compatibility examples, not
arithmetic counterexamples. The final example at \(p=37\) is an
arithmetic counterexample only to reversing the reflection implication.

\begin{proposition}[K17-GAMMA-NORMAL-CLOSURE]
\label{prop:gamma-normal-closure}
In the selected-free-factor construction, write
\(M=H^{\rm ab}\), \(T=T^-\oplus T^+\), and \(Q=M/T\).
Let \(s:Q\to M\) be its continuous \(\Delta\)-equivariant
\(\mathbb Z_p\)-section. Let \(T_I\) be the sum of the selected
even character components and \(T_J^+\) the sum of the omitted
even components. Let \(N\) be the smallest closed subgroup of
\(H\) containing \(H_I\) and normal under both \(H\) and
\(\Gamma\). Define
\[
 d_s(q)=\gamma s(q)-s(\gamma q)\in T^-,\qquad
 C_s=\overline{\Lambda\langle d_s(q):q\in Q\rangle}.
 \tag{GN.1}
\]
Then
\[
 (H/N)^{\rm ab}\cong (T^-/C_s)\oplus T_J^+
 \tag{GN.2}
\]
as \(\Lambda\)-modules with \(\Delta\)-action. In particular,
\[
 N=H\quad\Longleftrightarrow\quad C_s=T^-\text{ and }T_J^+=0.
 \tag{GN.3}
\]
\end{proposition}

\begin{proof}
The defect belongs to \(T\) because its image in \(Q\) is zero.
It is odd because \(s\) preserves \(\Delta\)-characters,
\(Q\) is odd, and \(\Gamma\) commutes with \(\Delta\).
Normal conjugation by \(H\) has no effect on \(M\). Thus the
image of \(N\) is the closed \(\Lambda\)-submodule generated
by \(s(Q)+T_I\). It contains all defects and hence contains
\[
 S=s(Q)+T_I+C_s.
\]
Conversely, this compact sum is closed and \(\Gamma\)-stable:
\(\gamma s(q)=s(\gamma q)+d_s(q)\), while \(T_I,C_s\) are
stable. Applying the same identity to \(\gamma^{-1}q\) proves
inverse stability, since \(\gamma^{-1}\in\Lambda\). Consequently
the image is exactly \(S\). Inclusion of the remaining torsion
modules gives the equivariant isomorphism
\(M/S\cong (T^-/C_s)\oplus T_J^+\), proving (GN.2).
Any nontrivial pro-\(p\) group has a nontrivial finite \(p\)-group
quotient, which has nontrivial abelianization. Thus a pro-\(p\)
quotient with zero abelianization is trivial, proving (GN.3).
No inference that finite abelianization implies a finite group is used.
\end{proof}

\paragraph{The odd recovery depends on the section.}
Consider the split module
\(M=\Lambda X\oplus\mathbb Z_pY\), where
\(\gamma Y=\alpha Y\), \(\alpha\in1+p\mathbb Z_p\), and
both summands have the same odd \(\Delta\)-character. Put
\(T=\gamma-1\). The section \(s_0(f)=fX\) has \(C_{s_0}=0\).
The equally valid continuous \(\mathbb Z_p\)-section
\[
 s_1(f)=fX+[T]f\,Y,
 \qquad [T]f=\text{the coefficient of }T\text{ in }f(T),
\]
satisfies
\[
 \gamma s_1(1)-s_1(\gamma)=-Y.
 \tag{GN.4}
\]
Hence \(C_{s_1}=\mathbb Z_pY\), integrally. Omitted even
components remain a section-independent obstruction in (GN.2),
whereas omitted odd components can be recovered by this defect
even when the original module extension splits over \(\Lambda\).

\begin{proposition}[K17-GAMMA-CLOSURE-COUNTERMODEL]
\label{prop:gamma-closure-countermodel}
Faithfulness on a selected free factor does not in general pass to
its \(\Gamma\)-stable closed normal closure. This failure occurs
in a formal model with all the free-group, module, exact-character,
equivariance and dense free-prounipotent-target hypotheses of the
cubic criterion, excluding its nonvanishing hypotheses. In that
model every allowed free factor generated by the \(Q\)-basis is
faithful, and one such factor has \(\Gamma\)-normal closure all
of \(H\), while the full representation has a nontrivial kernel.
\end{proposition}

\begin{proof}
\par\medskip\noindent\textbf{An elementary example.}
First let \(H\) be free pro-\(p\) on \(x,y\), with \(p\) odd.
The formulas
\[
 \phi_t(x)=x,\qquad \phi_t(y)=x^{t/2}y x^{t/2}
 \quad(t\in\mathbb Z_p)
\]
define a continuous \(\Gamma=\mathbb Z_p\)-action:
\(\phi_s\phi_t=\phi_{s+t}\) and \(\phi_{-t}=\phi_t^{-1}\).
The involution taking both generators to their inverses commutes
with it; both orders take \(y\) to
\(x^{-t/2}y^{-1}x^{-t/2}\). Abelianization
\(H\to\mathbb Z_p^2\) is equivariant and injective on
\(\langle y\rangle\). Its \(\Gamma\)-stable normal closure is
all of \(H\), since killing \(y\) and \(\phi_1(y)\) also kills
\(x\). But \([x,y]\) is a nontrivial kernel element: its first
Magnus term is \(XY-YX\). This already disproves inheritance of
faithfulness by normal closure.

\par\medskip\noindent\textbf{The free-orbit model.}
Fix a topological generator \(u\) of \(1+p\mathbb Z_p\), an
odd residue \(a_0\) modulo \(p-1\), and
\(\alpha\in1+p\mathbb Z_p\). Let \(P_0\) be free pro-\(p\)
on \(\gamma\) and \(x_a\), one for each odd residue \(a\).
Let \(F\) be the kernel of \(P_0\to\Gamma\) which sends
\(\gamma\) to a topological generator and every \(x_a\) to one.
It is free pro-\(p\) by the pro-\(p\) Schreier theorem. The same
finite-partition construction used in the moment argument identifies
its compact free generating orbits as
\[
 x_{a,t}=\gamma^t x_a\gamma^{-t},\qquad t\in\mathbb Z_p.
\]
Its abelianization is \(\Lambda\)-free on the \(x_a\).
Indeed, for any compact abelian pro-\(p\) \(\Gamma\)-module \(B\), arbitrary
images \(b_a\in B\) define a homomorphism
\(P_0\to B\rtimes\Gamma\) by the free universal property.
Restriction to \(F\) gives exactly the arbitrary continuous
\(\Gamma\)-equivariant homomorphisms \(F^{\rm ab}\to B\).
This is the universal property of the finite sum of copies of
\(\Lambda\).

Set \(H=F*_p\langle y\rangle\), adjoining an infinite procyclic
free factor. Define
\[
 \gamma^s x_{a,t}=x_{a,t+s},\qquad
 \gamma^s y=y^{\alpha^s},
 \qquad
 \delta x_{a,t}=x_{a,t}^{\omega(\delta)^a},\qquad
 \delta y=y^{\omega(\delta)^{a_0}}.
 \tag{GN.5}
\]
The \(\Delta\)-action is induced on \(P_0\) by fixing
\(\gamma\) and taking the indicated unit powers of its other free
generators. Its finite order and its commutation with \(\Gamma\)
follow on those generators. The formulas on the \(y\)-factor have
the same properties. These are continuous actions, and
\[
 M=\bigoplus_{a\text{ odd}}\Lambda X_a\oplus\mathbb Z_pY,
 \qquad Q=\bigoplus_{a\text{ odd}}\Lambda X_a,
 \qquad T^-_{a_0}=\mathbb Z_p(\alpha),\quad T^+=0.
 \tag{GN.6}
\]

Exact character lifts, including lifts inside arbitrarily small
stable neighborhoods, are available here. The proof of
\cite[Lemma~2.1, pp.~3--4]{Sharifi2002} uses only a pro-\(p\)
group with the coprime cyclic action: its iterated character
projection preserves the abelian character projection, remains in
each stable open subgroup, and converges to the exact eigenlift.
Choose lifts of a convergent product basis of the graph section
in (GN.4), and use the zero graph in the other characters.
Together with \(y\) these give an actual compact free basis.
As in the earlier basis lemma, the map from the free group on
that convergent basis is Frattini-surjective, splits by freeness
of \(H\), and the split image is again Frattini-surjective in
the source. This proves an isomorphism without invoking
finite-rank Hopficity. The subgroup \(H_s\) on the graph-section
basis alone is therefore a closed free factor.

\par\medskip\noindent\textbf{The dense free target and the faithful factor.}
Let \(U\) be the rational free prounipotent group on one generator
\(s_m\) in each odd arithmetic weight \(m\ge3\). Give it the
actions \(\gamma s_m=u^m s_m\) and
\(\delta s_m=\omega(\delta)^m s_m\). Define
\[
 \rho(y)=1,\qquad
 \rho(x_{a,t})=
  \exp\left(\sum_{\substack{m\ge3\text{ odd}\\m\equiv a\ (p-1)}}
                 u^{mt}s_m\right).
 \tag{GN.7}
\]
This gives a continuous homomorphism. In a finite weight quotient
the sums are finite and continuous in \(t\). In a faithful upper
triangular matrix realization, the compact generator orbits have
uniformly bounded entries. An entry of any word is a sum of
products of at most the fixed nilpotence length of off-diagonal
entries. Thus all word entries have one uniform denominator;
inverses have the same property by their finite geometric series.
The generated closure is compact and pro-\(p\), as is seen from
the successive additive upper-diagonal quotients. The
\(\Gamma\)-action preserves it. Apply the free pro-\(p\)
universal property in this compact group semidirect \(\Gamma\),
then restrict from \(P_0\) to \(F\); send \(y\) to one.
The compatible finite weight maps define (GN.7). Its formulas
show exact \(\Gamma\times\Delta\)-equivariance.

The image is Zariski dense in each finite weight quotient. Within
each character, the images at \(t=0,1,\ldots\) give a Vandermonde
matrix \(((u^m)^t)\) on the finitely many abelian generators.
Its determinant is nonzero because the \(u^m\) are distinct.
The Lie algebra of the image closure therefore surjects onto the
target abelianization. A subalgebra of a finite nilpotent Lie algebra
which surjects onto the abelianization is the whole algebra:
induct on the lower central series and replace each entry of a
commutator by a lift from that subalgebra. This proves density.

Every allowed factor generated by the \(Q\)-basis is faithful
under \(\rho\). To see this directly, a nonidentity element of
such a factor has a nonzero first Magnus coefficient in a completed
tensor power of \(Q\). In the orbit coordinates its finitely many
character components are bounded multivariable power series.
The leading degree-one change from its chosen basis to \(Q\)
is an isomorphism; adding torsion or higher-degree terms to the
basis lifts does not change that first coefficient. The target
coefficient at matching odd weights \(m_1,\ldots,m_r\) evaluates
it at \(T_i=u^{m_i}-1\). Higher source degrees cannot contribute
to that first target length. If the element were in the kernel,
all these evaluations would vanish. In each residue class the
allowed positive weights contain infinite sequences with an
interior \(p\)-adic accumulation point. Successive one-variable
analytic identity theorems make the original bounded power
series zero, a contradiction.

For the graph section (GN.4), \(C_s=T^-\). There is no omitted
even component, so (GN.3) gives \(N=H\). Nevertheless
\(y\ne1\) lies in the kernel of \(\rho\), and its actual
companion error is
\[
 \gamma(y)(y^\alpha)^{-1}=1.
\]
Its cubic numerator is identically zero. This is the asserted
formal model: it retains all the stated algebraic structure and
density, but has no arithmetic nonvanishing input. It is not
claimed to be a Galois realization.
\end{proof}

\paragraph{K17-REFLECTION-NO-ITERATION.}
The established reflection direction is
\[
 A^{[1-a]}=0\ \Longrightarrow\ A^{[a]}=0
 \qquad(a\text{ even}),
 \tag{GN.8}
\]
not its reverse. A Bernoulli unit supplies both sides for its
particular character; repeating the involution \(a\mapsto1-a\)
then adds nothing. The duality with \(M^+\) changes the ambient
module and includes its Tate twist; applying the inverse
identification returns the original class component.
Coleman/Wake--Wang-Erickson freeness of \(M^{[1-a]}\) makes its
unramified quotient cyclic, not zero.

The reverse of (GN.8) fails arithmetically at \(p=37\).
Verified Vandiver gives \(A^{[32]}=0\), whereas
\[
 B_{32}=-\frac{37\cdot683\cdot305065927}{510}
\]
and \cite[Main Theorem~1.1, p.~151]{Ribet1976} give
\(A^{[1-32]}=A^{[5]}\ne0\). The Vandiver input is within the
verified range of \cite{HartHarveyOng2017}; the displayed rational
Bernoulli number and its factors were independently checked in
the positive-component certificate. This is a counterexample to
reversing reflection, not to the kernel conjecture.

If \(E\) is the set of known vanishing even components, the
selected-pair theorem uses \(E\cap(2-E)\); it does not enlarge
\(E\). At the level of the present displayed integer windows,
\[
 E=\{2j:-12\le j\le16\},\qquad
 E\cap(2-E)=\{2j:-12\le j\le13\},
\]
which is precisely \(\{-2j,2j+2:0\le j\le12\}\).
At small primes the corresponding sets are taken modulo \(p-1\),
with their actual residue collisions. The growing Bernoulli window
deletes further components where its units have been proved;
iteration alone supplies neither new Bernoulli units nor missing
negative characters. Thus neither normal closure nor repeated
reflection in these forms closes the remaining arithmetic kernel
gate.
